% Standalone Overleaf document; pdfLaTeX; no external figures or tables.
% Scientific execution protocol proposed after the completed Stage-3 readout.
% This file does not record a Stage-4 model run.
\documentclass[11pt,a4paper]{article}
\usepackage[margin=2.1cm]{geometry}
\usepackage[T1]{fontenc}
\usepackage[utf8]{inputenc}
\usepackage{amsmath,amssymb,booktabs,longtable,array,enumitem,graphicx}
\usepackage[hidelinks]{hyperref}
\setlength{\parindent}{0pt}
\setlength{\parskip}{5pt}
\setlength{\emergencystretch}{2em}
\setlist{nosep,leftmargin=1.5em}
\newcommand{\file}[1]{\nolinkurl{#1}}
\newcommand{\Rtest}{R_{59}}
\newcommand{\Rall}{R_{61}}
\newcommand{\E}{\mathbb{E}}
\title{Stage 4: Causal Groups, Communication and RI Participation\\
\large A bounded execution protocol after the complete Stage-3 results}
\author{Phase A: Single-hop SIH Audit}
\date{Protocol v2: route-first, reduced RI audit}
\begin{document}
\maketitle

\section{Target claim and strategy}
Identify head groups that causally support the model's in-context mother-of
behavior, identify specific communication routes between them, and determine which
RI-selected candidates participate, including conditional and opposing contributions.
The semantic claim is task-specific: selecting the mother bound to the queried child
under the tested fact assignments and order changes; query-change claims require the optional query panel. It is not a claim of a general semantic
faculty, relation specificity across domains, or the discovery of every valid circuit.

Start with a finite list of exact path tests; construct group interventions from the measured routes. Use an
approximate screen only if a stated expansion trigger occurs. A faithful head subset
with tested communication is an acceptable result; a partially explained mechanism
or failure of the faithfulness criterion must also be reported. Do not force a single
chain or a single globally minimal circuit.

\textbf{Status.} This is a proposed scientific execution protocol. The accompanying
\file{results/stage4_design_v2/manifest_proposal.json} fixes seed rosters, reference
cohorts, source hashes and decision defaults. Its preparation involved saved-data
analysis only. No Stage-4 effect has been measured. Before GPU execution, the
implementation must materialize the specified configurations and pass Section~
\ref{sec:gates}; runtime and code hashes then complete the execution manifest.

\subsection{Evidence that determines this design}
The completed Stage-3 tables, not the older handoff's superseded interpretation,
give the following starting points.
\begin{itemize}
\item L17H1 has Scope-P importance $+5.701$ but Scope-F $-0.056$. Its single-site
effects at query-fact ``is'' and period are $+1.769$ and $+1.377$. The controlled
identity readout at ``is'' changes by $-0.458$ after replacing the head and by
$-0.521$ after source corruption, with the head shift negative in all 20 families.
This motivates a content route; it neither identifies a receiver nor proves a binding.
\item L15H25 has $P=+0.679$, with $+0.429$ at child-last. Its readout is inconclusive
because the full source-swap control barely changes the mother contrast. Retain this
writer and test its communication; do not require another successful readout first.
\item Twenty of 25 strong heads have $F/P$ near one. L17H3 is mixed-position
($P=-0.546$, $F=-0.402$). On the common pairs its non-query mother spans have summed
mean importance $-0.053$; no individual early role is a strong established write site.
\item Signed averages conceal large family interactions. For example, L27H6's
attention/value interaction averages $+0.205$ but its mean absolute family effect is
1.254. Group and path effects must therefore be measured, not added from singles.
\item \textbf{Eight}, not only five, of the 25 strong heads are test-only RI
candidates: five positive and three negative. L19H16, L23H15 and L26H23 are in RI.
The mediation roster is therefore not disjoint from RI.
\item Re-reading all 1,024 exact head effects on the original 40 pairs finds two
outside the Stage-3 inventory above $0.3$: L13H18 ($+0.387$) and L24H19 ($+0.390$).
Their all-discovery effects and positions are not established. They receive a small
mandatory coverage check rather than being silently excluded from circuit search.
\end{itemize}

\subsection{Changes to the practical proposal}
K/V for a fact-to-answer route are replaced at the \emph{fact/source positions};
Q is at the receiver's query position, usually the colon. A same-layer head cannot
read another head's current-layer output. L26H31 cannot change its attention in
response to ablating downstream L27H6. A group ablation is not an upper bound on
mediated importance. Residual caches required here are not present in the saved
Stage-3 projection tables and must be collected. A fixed final-normalization scale
alone does not handle OLMo's shared attention-output normalization. Random-edge
percentiles and an 80\% sign rule do not establish significance or universal membership.

\section{Populations, metrics and decision rules}
\label{sec:metrics}
Use the pinned \file{allenai/OLMo-2-1124-7B} revision
\file{7df9a82518afdecae4e8c026b27adccc8c1f0032}, tokenizer and deterministic eval
runtime. Preserve the original split: 89 discovery families, two orders each;
the same 20 common families/40 pairs for initial exact tests; 87 held-out families
opened once after the circuit and evaluation suite are frozen. The manifest lists
the common family IDs. The unused 69 discovery families provide an expansion check,
not a pristine confirmatory sample after previous Stage-2 use.

For each pair, use its clean and corrupted gold answers at their first divergent
token, with their shared answer prefix teacher-forced:
\[
 M(x)=z(y_c\mid x,c)-z(y_r\mid x,c),\quad
 I_T=M(x_c)-M(x_c;T_{r\to c}),\quad
 J_T=M(x_r;T_{c\to r})-M(x_r).
\]
Positive $I$ means noising moves the answer toward corruption; positive $J$ means
restoration moves it back. Negative contributors remain eligible. Save both logits,
candidate probabilities and full-vocabulary top token, not just the difference.
Average orders within family and then families equally. Also retain per-order values,
absolute family effects, median, range and sign fractions. Never substitute
$|\E I|$ for $\E|I|$.

The mean full-model gap on the saved 178 pairs is 11.17046 logits (11.20006 on the
common 40). Report normalized effects as a ratio of means to the gap on the identical
population, alongside raw effects. Do not average per-pair ratios. If the denominator
is below 0.5 logits in a subgroup, suppress its ratio. The four shared-prefix pairs
are families 168 and 174, both orders. During route discovery, patch only original
prompt positions and let the prefix recompute. Include these pairs in full discovery;
report a no-shared-prefix sensitivity rather than dropping them silently.

\textbf{Operational retention, fixed before new effects.} The practical mean-effect
threshold is $\tau=0.10$ logits, a design choice, not a significance threshold.
A coherent route/member candidate has $|\bar I|\ge\tau$ and the selected mean sign in
at least 70\% of families. Also retain a heterogeneous candidate if
$\E_f|I_f|\ge\tau$ and at least 20\% of families have $|I_f|\ge\tau$; its label must
remain heterogeneous until tested conditionally. Do not reject such a candidate merely
because its signed mean cancels. The only prespecified context strata are query-fact
before versus after the competing fact, and corruption axis; require at least ten
families per reported stratum. No post-hoc token/name-based subgroup fishing.

Use 20,000 whole-family bootstrap draws, seed 20260923, preserving all paired
conditions/orders. Discovery intervals are descriptive after selection. On validation,
freeze each claimed head's background, route, site and sign in advance. Report ordinary
95\% family intervals and simultaneous 95\% intervals for the vector of claimed
signed mean effects: use the 95th percentile of the bootstrap maximum absolute
standardized mean deviation, with the original family standard errors held fixed.
Flag zero-variance columns explicitly. A confirmed directional contribution must have
the frozen sign, magnitude at least $\tau$, and a simultaneous interval excluding
zero. An inconclusive interval is not evidence of absence. Nonlinear faithfulness
ratios are bootstrapped jointly with their denominators.

\section{Head rosters and positional masks}
\label{sec:rosters}
\begin{longtable}{p{.22\textwidth}p{.72\textwidth}}
\toprule Roster & Members and use \\\midrule
Primary writers & L17H1; L15H25. \\
Secondary early/mixed & L13H10, L14H26, L17H3. \\
Positive group $A$ & L27H6, L21H18, L22H5, L21H6, at the colon. \\
Negative group $B$ & L20H1, L19H16, L26H23, at the colon; causal sign does not imply a mover mechanism. \\
Mother-attending priority & L18H18, L18H19, L19H22, L21H6, L21H18, L22H5, L27H6. \\
Child-attending priority & L16H1, L16H21. \\
Additional colon sources & L17H24, L25H17, L21H23, L30H13, L17H17; self attention is not exclusive fact avoidance. \\
Additional negatives & L23H15, L30H18, L18H24, L17H3, L13H10. \\
RI strong positive & L16H1, L16H21, L17H24, L18H19, L22H5. \\
RI strong negative & L19H16, L23H15, L26H23. \\
Backup candidate & L26H31, also in RI. \\
Early minor panel & L6H24, L7H1, L8H15, L9H16, L12H17. \\
Outside-inventory reserve & L13H18, L24H19; exact 40-pair Scope-P evidence only. \\
References & The unchanged 25 Stage-2 random heads; four fixed 23-head layer-matched reference cohorts for the reduced RI residual group. \\
\bottomrule
\end{longtable}

$\Rtest$ is the original 59-head test-only RI union; the two historical extras
L3H11 and L9H22 are not part of the mandatory v2 audit. Define $R_{31}$ by retaining
an original candidate if its best eligible Stage-1 selection rank is at most three
in any original scope/anchor/metric, OR if
$\max(|P_h|,|F_h|)\ge0.10$ logits in the saved full-discovery effects.
This keeps 23 by RI rank and eight additional causal candidates, dropping 28.
It is a prioritization rule, not a significance or semanticity threshold.
All eight strong RI heads and L26H31 remain. The CSV records every decision.
The residual group $R_{23}$ comprises the 31 minus the five strong positive and
three strong negative heads; it is not the set of the 23 rank-selected heads.

The initial route pool is the 25 strong Stage-3 heads, including five source
candidates L13H10, L14H26, L15H25, L17H1 and L17H3. Source and receiver roles can
overlap; they are not exclusive classes. At a fact site, use later heads as
receivers; at the colon, test eligible later heads and downstream MLPs/output.
There are 21 heads with $|F|\ge0.3$ (14 positive, seven negative), and four further
strong Scope-P heads. In particular 16 of the 21 are later than layer 17.
The two coverage candidates and supported G2/G3 contributors can extend this pool.
The old four-positive/three-negative roster is a background-selection reference,
not a fixed coalition to exhaustively enumerate.

\textbf{Source sites.} L17H1: query-mother token span, ``is'', period, each separately;
also their union. L15H25: child-last, with ``is'' as a low-effect position control.
For writer-group restoration use two masks: these narrow sites, and the entire query
fact sentence. Include same-role positions in the other answer-side fact as controls.
For L13H10/L14H26, start with the entire query fact; for L17H3, separate colon,
non-query mother spans, and remaining earlier test positions. Derive indices from
token offsets and fact annotations, assert the decoded text, and keep all masks in
the manifest. Never treat the sum of site effects as a joint effect.

\textbf{No answer-selection oracle in circuit evaluation.} Query-relevant sites can
be used to localize interventions. A deployed circuit mask, however, must retain the
corresponding fact roles in \emph{all four facts}, not just the fact known to contain
the answer. Otherwise the experiment performs part of the semantic selection outside
the model. Question-token and colon masks may depend on syntax, not answer identity.

\section{Intervention implementation and gates}
\label{sec:gates}
\subsection{Two baselines}
Use paired-donor replacement for exact paths, G1--G3 and writer restoration. Use
discovery-mean replacement for circuit retention/removal, and as a sensitivity check for retained G1--G3 findings. Circuit conclusions also receive paired-donor sensitivity. Replace a head's
$z_h=A_hV_h$ slice at the input to \texttt{o\_proj}, preserving all other slices.
Replace MLP output at its residual branch boundary \emph{after} the MLP's shared
post-feedforward normalization when freezing that branch. Keep all normalization
operations in their declared live/frozen positions; do not normalize isolated heads.

Means are estimated from the 178 discovery pairs, balanced over family, order and
clean/corrupted variant. Match layer/head, test fact slot (textual slot 1--4), token
role and within-span bucket (single, first, interior, last; fixed template tokens
separately). Question names use the same token-role buckets. Do not condition on
query relevance, candidate identity or gold answer. Average events within each
family before families. For fewer than ten supporting families, drop fact slot,
then within-span bucket, then use the component's all-test-position mean, recording
the fallback. Discovery evaluation uses leave-one-family-out means; validation uses
the frozen full-discovery means. Demonstrations and shared answer-prefix sites stay
live in the primary scope, explicitly recorded below.

\subsection{Exact channel/path intervention}
\label{sec:path}
This is an OLMo-specific path intervention inspired by the controlled-path and
patching methodology of \cite{ioi,patch}, with its normalization boundaries specified
explicitly rather than assumed additive.
An attention path is $(u,S)\to(v,T,\gamma)\to$ answer, with
$\gamma\in\{Q,K,V,KV\}$. A fact-site source feeding a reader's colon output uses
$T=S$ for K/V and reader output row colon. Q at the colon has no direct same-position
residual route from a different fact token. Such a cross-position Q effect requires
a live intervening attention operation and must be named as a mediated path.

For each pair and source, cache clean/corrupted source $z$, branch outputs and actual
post-QK-normalization/RoPE Q,K plus actual V. Implement:
\begin{enumerate}
\item In a clean recipient, insert the corrupted source at $S$. For a
\emph{direct residual route}, freeze intervening post-normalization attention and
MLP residual-branch outputs to clean values. The source attention-output normalization
is live; its downstream same-layer MLP is frozen unless it is the receiver. Capture
the receiver's selected channel at $T$ before the receiver produces its output.
Receiver Q/K normalization and RoPE run normally before capture.
\item In a fresh clean run, insert only that hybrid channel slice at its matching
SDPA boundary. Keep other channels, rows and positions clean at the receiver;
subsequent model computation is live. Measure $I_{u\to v,\gamma}$. The source is
not patched in this final run. Reverse recipient/donor roles for restoration $J$.
\item For an MLP receiver, capture the changed MLP input at the same source position,
inject it in a fresh clean run, and let its output and downstream computation run.
The same-layer MLP after a source head is temporally eligible; a same-layer head is not.
\end{enumerate}

Freezing an intervening pre-normalization head sum while leaving its normalization
live would create an unintended extra path. Freeze the \emph{actual residual branch
increment} there. At the source itself the transmitted increment is
\[
 \delta_u=N_{\rm attn}\!\left(\sum_h o_h+o_u^r-o_u^c\right)
             -N_{\rm attn}\!\left(\sum_h o_h\right),
\]
which depends on the other clean heads. Group replacements use one joint
normalization; do not sum individually normalized increments. Direct-route results
are conditional on this clean reference. Reverse runs use their own corrupted
reference. Record the exact frozen/live branches with every route.

For mediated paths, unfreeze only the named intermediate receiver output/branch and
repeat the hybrid construction. A failure of the frozen direct route does not rule
out mediation. An additional all-mediators-live capture can locate total source
influence on a receiver, but must be labelled total influence rather than a direct edge.

\subsection{Required implementation gates}
Before scientific selection, test six fixed discovery pairs: the first two core
families in both orders and family 168 in both orders. Require unchanged-hook and
self-patch logit drift $\le10^{-3}$; saved Stage-2/3 baseline/patch replication
$\le0.05$; reconstructed SDPA output max error $\le0.02$; and combined QKV receiver
replacement versus ordinary receiver-output replacement readout drift $\le0.05$.
Test same-layer head-to-head and future-to-past null paths ($\le10^{-3}$), direct
cross-position residual nulls, mean-mask all-retained equivalence, finite values,
causal masks, and exact gold/prefix alignment. Test the path code on a tiny model
with an explicit manually isolated route before these real-model gates.

For endpoint/no-op gates, compare all measured answer logits, not only top-five
outputs or differences that could cancel. For joint group replacement, demonstrate
that source masks are applied simultaneously; restoring a downstream head from an
unmodified cache otherwise overwrites upstream influence. Stop on failed gates;
do not reinterpret a numerical failure as a negative finding or relax tolerances
after observing scientific effects.

\section{S4.1: route mapping before coalition selection}
\label{sec:routes}
\textbf{Coverage check.} Complete Scope P and F on all 178 pairs for L13H18 and
L24H19, reusing their existing 40-pair P values. If either meets the retention rule,
add it to the exact candidate roster. If $|F/P|<0.8$ with $|P|\ge0.1$, localize it
on the common pairs using six joint masks: all fact-mother spans, fact-child spans,
fact template words, fact punctuation, question tokens and final colon. Split a
retained fact-role mask by textual fact slot before path testing. These effects are
joint mask measurements, not sums of unmeasured single-position effects.


\textbf{Staged execution.} P tests nominate answer-relevant communication, not merely
activation changes. Start on the common 40 pairs. P1/P2/P4 map candidate sources;
P3 checks downstream continuation before interpreting a group as parallel branches.
Do not wait to exhaust all possible routes before testing a supported group.
All retained effects use the existing coherent/heterogeneous rule, including either sign.


\textbf{P1: L17H1 routes.} For each of its three source roles, test K and V separately
into every strong head in layers 18--31 (16 heads, including positive and negative
ones). The receiver output row is the colon and the injected K/V positions are the
source role. This is 96 attention-channel configurations on the common pairs.
Initially test same-layer and next-layer MLPs at these roles; later MLPs are conditional expansion candidates, not a mandatory 45-configuration sweep.
For any receiver where either channel is retained, measure joint KV and its
interaction; test a union of the three source roles rather than adding them.
L17H1 $\to$ L17H24 is a temporal null, not a candidate consumer. An effect at ``is''
does not imply that a head attending a mother-name token consumes it.

\textbf{P2: L15H25 routes.} At child-last test K and V separately into all 21
strong answer-position candidates later than layer 15, starting with L16H1/L16H21.
Use the colon output row. For retained receivers test joint KV; use ``is'' as the
low-effect site control on those receivers. Initially test MLPs 15 and 16 at child-last.
A failed readout does not exclude this source; unresolved strong influence triggers mediation.

\textbf{P3: colon outputs and mediation.} Sources are L16H1, L16H21, L17H24,
L18H18, L18H19, L19H16, L20H1, L23H15 and L25H17. Prioritize sources reached by P1/P2/P4 and retained G2/G3 contributors, then remaining listed sources. For each, test its colon output
into (i) every later strong head's Q at the colon and its joint KV at the colon,
(ii) its same-layer and next-layer MLP at the colon (more distant MLPs only on expansion), and (iii) a
residual-to-logits route. The latter freezes all later residual branches to clean
values but keeps the exact final normalization live; it measures the isolated
residual bypass under that reference, not a linear vocabulary projection. Split a
retained joint KV route into K and V. Do not first label these heads ``indirect''
from their projection gap; the route results decide what can be said.
On shared-prefix pairs, a colon-only residual bypass cannot directly reach the
later metric position when all intervening branches are frozen. Its null effect is
structural, not evidence against the source. Report those pairs separately and
require an explicitly live prefix-position attention mediator to test that route.

\textbf{P4: secondary sources.} Test L13H10 and L14H26 at the whole query sentence
against all later strong answer-position K/V receivers, prioritizing receivers retained for other sources, and their own-layer/next-layer MLPs.
L17H3's main source is the colon; apply the P3 recipe once, deduplicating overlapping configurations. Test its non-query mother spans and remaining
earlier positions as joint auxiliary masks, not a confidently localized early write.
Only if a source or route remains unexplained, the five early minor heads receive one common-pair source-group test each into the
retained receiver groups; a null result means no route detected in that panel, not
permission to claim those heads cannot participate elsewhere.

\textbf{Controls and extension.} Every route gets a self-donor check; impossible
routes are implementation controls, not a percentile reference. For the 24 largest
retained routes include a same-source wrong-role position and two seeded,
same-layer non-roster receiver heads; report intervention norms rather than assuming
random effects have matched magnitude. Both effect signs are eligible. Extend
retained routes to all 178 pairs and reverse recovery, including a separate summary
on the 69 additional families. Prioritize one route per retained RI head and per
effective writer before ranking further routes by absolute family effect. Cap the
full-discovery panel at 96 routes; overflow stays explicitly untested rather than
being declared absent.

\section{S4.2: route-supported groups and reduced RI audit}
\label{sec:groups}
\textbf{G1: construct groups from measured routes.} For each source/site, form the
set of recipients with retained exact answer effects. Conversely, for each
receiver/site/channel form the set of supported sources. Preserve layer order,
channel, position and effect sign; a chain is not a set of independent branches.
An unexplained strong reader remains a search candidate even if no initial writer
connects to it. Source labels never veto an incoming route.

On the common pairs, test groups with at least two members using simultaneous
paired-donor replacement at the mapped sites: each member alone, the whole group,
and the group minus each member. For a group $S$ of size $n$, this is at most
$2n+1$ nonempty configurations before deduplication and reuse, not $2^n$.
Save $I(S)-\sum_{h\in S}I(h)$ and
$\delta(h\mid S\setminus\{h\})=I(S)-I(S\setminus\{h\})$.
Leaving one member active tests its conditional contribution; call it sufficient
backup only if the remaining task performance also supports that interpretation.
Run the receiver-blocking test below for supported source-to-group patterns.

Bound selection to eight distinct groups initially. First take one outgoing group
per effective source in fixed order L13H10, L14H26, L15H25, L17H1, L17H3; then
incoming shared-receiver groups, ranked by the sum of mean absolute family route
effects, with layer/head/site/channel tie-breaking. That ranking is only a priority
score, never an estimate of joint importance. Deduplicate identical masks.
Use at most 128 distinct common-pair group configurations. If a group would exceed
the remaining budget, record it as untested rather than silently truncating members.
Only a retained joint nonadditivity or conditional effect triggers pairwise splitting,
within that same cap. Extend at most 16 informative group contrasts and their required
baselines to 178 pairs, ranking by mean absolute family contrast with deterministic
ID ties. Retained claims receive mean-replacement sensitivity; cancellations are
reported per family rather than hidden by the signed mean.

\textbf{Receiver blocking: a focused communication test.} For each supported source, choose its route-supported receiver group. Test joint source restoration only when sources share a supported receiver or a measured interaction motivates it. Keep head and MLP groups separately identified; do not combine the two primary writers by default.
Compute in corrupted recipients
\[
 J_W=M(x_r;W_c)-M(x_r),\qquad
 B_{W,V}=M(x_r;W_c)-M(x_r;W_c,V_r).
\]
Here $V_r$ clamps the receivers' \emph{outputs at their tested positions} to their
original corrupted values, preventing their response to restored writers. Also run
$M(x_r;V_r)$ as the self-clamp control. A positive $B$ shows that allowing these
receivers to respond is needed for part of the measured writer restoration in this
background. Reverse the experiment: noise $W$ in clean inputs and restore receivers
to clean values. Report raw effects, source/receiver joint effects and interactions;
do not call $B/J$ a mediated fraction unless the denominator is stable, and even then
label it an intervention ratio. Overlapping receiver groups need not add to one.
This test complements isolated paths with communication in the otherwise live model.


\textbf{G2: focused RI participation.} Use the 31 selected heads, not all 59/61.
Reuse existing Stage-2 single-head paired-donor Scope-P effects in the intact model,
checking exact model/input/site/metric agreement. New effects use the same full
original prompt mask and donor policy (identical-prefix sites have zero replacement
change), not a narrower colon mask. Recompute only missing or incompatible baselines.

Choose a weakened-reader background $K$ on the common pairs from prefixes of
(L21H18, L22H5, L21H6, L27H6), patched at the colon: choose the smallest prefix
leaving 30--80\% of the original gap, otherwise the closest fraction to 0.55 and
flag floor effects. Record $I(K)$ and, for every $h\in R_{31}$, measure
$I(K\cup\{h\})-I(K)$; if $h\in K$, use $K\setminus\{h\}$ as its background.
Apply every donor mask simultaneously. Compare with the saved intact-model single.
This requires 31 times 40 = 1,240 initial conditional intervention evaluations,
plus background controls; no duplicate intact-model single sweep is required.
Extend heads meeting the retention rule in their conditional effect or its change
from the intact effect to all 178 pairs. Always carry the eight established strong
RI candidates into route follow-up even if a conditional mean cancels.

Test the five strong positive, three strong negative and remaining 23 as three
fixed groups in intact and weakened backgrounds (incremental effect relative to
$K$, deduplicating overlaps). Add four fixed reference cohorts of 23 non-RI,
non-strong heads, matched exactly to the residual group's layer histogram, in both
backgrounds. The manifest materializes these sets with PCG64 seed 20260923; exclude
all original/historical RI and coverage reserves. These are descriptive controls,
not null significance thresholds. Start on 40 pairs; extend retained contrasts only.
No exhaustive subsets are taken. Mean replacement is a sensitivity test for retained
claims, with both intact and weakened comparators recomputed consistently.

\textbf{G3: L26H31 backup test.} This can run early alongside P mapping once gates
and donor baselines pass; it does not require a completed route map or G1.
Measure its incremental effect in backgrounds $\varnothing$, $\{$L21H18$\}$,
$\{$L21H18,L22H5,L21H6$\}$ and that set plus L27H6, at the colon. Reuse compatible
G2 conditions, otherwise keep the different source masks explicit. Capture its
attention/output before its intervention; include reverse restoration for retained
contrasts. Only upstream removals can change its attention: L27H6 removal alone is
an attention-null control, though it can change the answer effect of L26H31.
Call it backup only with a supported conditional contribution, not from attention
alone. Start on 40 pairs and extend retained contrasts to 178.

\textbf{Scheduling.} The reading order is P mapping, G1, G2, G3. Execution can
interleave G2/G3 with P mapping; their supported contributors feed back into routes.
S4.3 is a triggered return loop, not a mandatory barrier after all G tests.

\section{S4.3: bounded expansion only where needed}
\label{sec:expansion}
Trigger expansion if a writer has a substantial total effect/restoration but no
supported consumer, an RI conditional contributor lacks a route, or the candidate
head mechanism fails discovery faithfulness/completeness. Do not expand merely to
obtain a desired story. Run at most two rounds, adding at most 12 heads and four MLPs
per round, with untested boundaries recorded.
Limit each round to 256 new exact configurations. Prioritize a route for every
unexplained primary writer and retained RI contributor before sorting remaining
candidates by mean absolute family effect.

\textbf{First expansion: exact local mediation.} For an unresolved writer, allow
its own-layer and next-layer MLP branches live in the hybrid construction, one and
then both, and recapture the previously tested receiver channels. For an unresolved
colon source, use its same-layer MLP and the next two blocks as mediator candidates.
Test actual ordered routes; reintroducing a whole block establishes a block-mediated
route, not every head inside it. Split an effective block into its 32 head outputs
and MLP to select the next route tests. Preserve same-layer attention causality.

\textbf{Second expansion: broad candidate search.} Search all 1,024 head outputs
and 32 MLP branches, without using RI membership as a filter. Use actual
head-output attribution under the current full or weakened/circuit background,
stratified by fact spans, question and colon; the score is
$(z^{\rm donor}-z^{\rm recipient})^\top\nabla_z M$ through the real model, including
normalization. Rank signed means and mean absolute family scores separately.
This is a node screen to nominate exact tests, not an edge claim. Saved Stage-2
all-head scores are an additional priority list, not calibration for this background.

Calibrate on 48 candidates fixed after ranking but before their exact effects:
12 highest absolute mean, 12 highest mean absolute excluding those, 12 middle,
12 near-zero/random, spanning layers/sites. Accept prioritization only with Spearman
$\ge0.8$ and sign agreement $\ge0.9$ among exact effects above $\tau$; report exact
strong effects in the low-score strata. If fewer than ten above-threshold examples
exist, sign calibration is insufficient. A failed screen falls back to exact tests
of the 105-head inventory plus eight four-layer outside-inventory blocks, splitting
effective or cancelling blocks under the remaining budget. This fallback bounds the
search; null blocks do not prove all their heads are irrelevant.

After adding candidates, return to exact channel paths and conditional contributions.
An adapted edge EAP/EAP-IG implementation is optional only if the unresolved
receiver-specific edge list exceeds 500 configurations. It must use the exact
intervention Jacobian and fresh exact calibration above; ordinary EAP's
two-forward/one-backward cost is not promised for OLMo's nonlinear branch graph
\cite{eap,faith}. No approximate score alone establishes or excludes membership.

\section{S4.4: test the task-level relational behavior}
\label{sec:semantic}
\textbf{Optional four-condition panel, not an entry gate.} Existing intact-model
base, corrupted and 40 changed-query measurements already support task capability.
Do not repeat a capability survey or expand it merely for new pairs. Enable this
panel only if a specific unresolved mechanism question concerns joint fact/query
selection; freeze that decision before held-out validation. Reuse compatible saved
baselines and measure missing intact outputs only for actual intervention comparisons.
No optional-panel result is required to start P/G experiments.

\textbf{If enabled:} For each common family/order, cross original/swapped
mother assignments with original/alternative answer-side queried child:
$x_{00},x_{10},x_{01},x_{11}$. Keep demonstrations fixed and reconstruct gold answers
from the fact table. If the two candidate mothers are $a,b$, the correct answers
are $a,b,b,a$. With a fixed $a-b$ answer-token sign, the full model should therefore
reverse under either factor and return under both. Measure
\[
 B_D=\E_{f,o}\,[M_D(x_{00})-M_D(x_{10})-M_D(x_{01})+M_D(x_{11})]/4.
\]
This is a crossed behavioral contrast, not a neuron semanticity score. Report all four
condition accuracies/margins, not only $B_D$. Preserve model errors; the existing
38/40 query-change success does not guarantee the new crossed condition works.

Run the full model, the retained head mechanism, its RI-member removal, its
non-RI-member removal, and the writer-restored/receiver-blocked mechanisms selected
above (at most six configurations). Apply mean baselines symmetrically. Compare
which intervention disrupts fact-assignment following, query following, or both.
For route-level contrasts use the four adjacent edges of this square with their
own recipient gold sign and first-divergent prefix, not the fact-swap sign blindly.

\textbf{Important causal constraint.} Under query-only corruption, all fact-site
activations preceding the changed question are identical by causality. A fact-site
writer donor patch on that axis is therefore a deliberate null. Test query selection
through receiver Q at the colon and through question/colon source heads identified
in P3/expansion; do not classify a writer as lexical merely because its fact-site
effect is zero on the query axis. The meaningful mechanism can combine a
fact-dependent V route with a query-dependent Q route.

Use tokenizer offsets to map unchanged tokens and annotated first/last anchors.
Query names may differ in length: never reuse Stage-2's equal-length/six-token-change
assertion for this new axis. Recompute causal masks, RoPE positions and prefixes;
unmatched interior spans are unavailable for that route, not padded donor vectors.
Use only original positions for route claims and disclose unsupported alignment.
Do not discard entire families from behavioral or circuit evaluation.

\textbf{Order robustness.} Evaluate the same frozen mechanisms on the existing
reorder variants with role-mapped masks. Report absolute family changes and
condition-specific performance, not invariance inferred from a near-zero signed mean.
Extend selected order-robustness contrasts to all 89 discovery families before validation. Extend the four-condition panel only if explicitly enabled for a mechanistic question.
This tests in-context binding/selection in the existing relation. Location-world
transfer, parent coverage and a new observational RI are excluded.

\section{S4.5: assemble and evaluate a head-level mechanism}
\label{sec:circuit}
\textbf{Primary achievable claim.} A position/role-defined attention-head mechanism
with a tested communication map, conditional on live MLPs, demonstrations, embeddings,
normalization and answer-prefix computation. It is not a fully isolated model circuit
or an edge-minimal graph. This scope makes the central RI question answerable without
first solving every nonlinear MLP/normalization interaction in OLMo.

Start $C_0$ with the 25 strong Stage-3 heads plus new exact/conditional contributors.
Initially retain each included head at all test-block positions; all other head
slices there receive discovery means. This broad mask avoids losing an unlocalized
source before testing it. Test finer role masks only by direct removal, broadening
query-specific sites to all fact slots. MLPs stay live. Also evaluate the all-head-
ablated baseline with this same background. If its gap already approaches the
80\% criterion, report that the criterion does not isolate the head mechanism and
include MLP interventions before making a self-contained circuit claim.

For each family let $g_f(D)$ be its mean clean-minus-corrupted gap across orders.
Evaluate
\[
 F(C)=\frac{\E_f g_f(C)}{\E_f g_f(\mathrm{full})},\qquad
 L(C)=\frac{\E_f|g_f(C)-g_f(\mathrm{full})|}
 {\E_f|g_f(\mathrm{full})|}.
\]
The inherited criterion remains $F\ge0.80$. To call the result a useful faithful
mechanism, additionally require $F\le1.20$, $L\le0.20$ and candidate-accuracy loss
no greater than five percentage points separately on clean and corrupted inputs.
Report full-vocabulary first-token accuracy and per-family behavior as secondary
checks. These added guards are prospective safeguards against overshoot/cancellation,
not retroactive changes to Stage 3. On the reorder panel, and on the four-condition panel only if enabled, require
the same five-point condition-wise accuracy guard relative to the full model;
report $B_C/B_{\rm full}$ only with a stable denominator. Negative-head removal that
improves one metric does not by itself make a more faithful explanation
\cite{faith,robustness}.

\textbf{Reduction.} Use the common pairs to greedily remove the head/role with the
smallest measured increase in $L$, breaking ties by layer/head/role, only when all
guards still pass. Re-evaluate contributions after every accepted removal. Limit
to two full passes, then evaluate the resulting mask on all 178 pairs. If it fails,
restore removed elements in reverse order until the guards pass, recording both
curves. Do not choose another metric to salvage a failure. There is no claim of
global or unique minimality.

\textbf{Completeness and backup.} In the full model and $C$, remove the same
predefined groups: retained writers, child readers, positive readers, negative
contributors, RI members, and the selected weakened-reader background. Compare
$g_f(\mathrm{full}\setminus K)$ with $g_f(C\setminus K)$ and both accuracies.
A normalized absolute gap discrepancy above 0.20 or more than five points of
condition-wise accuracy discrepancy triggers the bounded expansion/addback step.
In addition test single-head addback for retained or conditionally supported $R_{31}$ heads absent from $C$, plus up to
16 non-RI candidates ranked by prior exact conditional evidence. Absence from one
compact faithful subset does not prove absence from the full model's alternative
mechanisms. Whole-$C$ removal from the full model is useful evidence, but its failure
to destroy behavior may reflect backups; joint group importance is not an upper bound.

Repeat final retained/removal configurations with paired-other-condition donor
replacement, applied symmetrically to clean and corrupted inputs. Report ablation-
dependent results honestly. If the mean and donor stories conflict, call membership
baseline-sensitive until the conflict is resolved \cite{best,robustness}.

\textbf{Optional stronger circuit claim.} Only if needed after the head-level result,
include and ablate MLP nodes, prefix positions and excluded paths. A claim of an
edge-defined circuit requires actual excluded-edge masking with shared normalization
explicitly retained in the computation graph; node ablation plus a path drawing does
not establish edge faithfulness. This is a separately budgeted extension, not an
implicit claim attached to the primary head-subset score.

\section{Membership and the SIH/RI conclusion}
\label{sec:membership}
A \emph{supported head member} needs both (a) an exactly tested route to an answer-
affecting receiver or residual-to-logits path, possibly through named mediators,
and (b) a conditional removal/restoration contribution of at least $\tau$ in a
declared full-model or circuit background. The latter may be negative. Specify the
site and whether the effect is coherent or context-specific. Passing a group test
alone does not establish every member; raw vocabulary alignment alone does not
establish communication. A conditional contributor whose route is unresolved is
reported separately and is not erased from the RI audit.

Report the original 59-head ledger, marking the 28 dropped heads ``not selected
for the v2 targeted audit'', not absent. For the 31, report saved baseline effect,
conditional effect, site, route, sign, removal/addback, sensitivity and validation
status. Historical extras are historical only unless independently rediscovered.
Report $|R_{31}\cap C_H|/31$, $|\Rtest\cap C_H|/|C_H|$, and the count of supported
conditional members outside $C$. A count divided by 59 is observed coverage, not a
full-list exclusion or precision estimate. Selection partly used causal effects;
do not present enrichment of $R_{31}$ as independent validation of RI ranking.

The possible conclusions are: RI candidates participate in these mechanisms;
participation is selective or includes opposing/backup functions; no contribution
was detected under the bounded tested scope; or coverage/precision is insufficient.
Even a clear positive result establishes participation of the implemented RI
candidates, not that RI is a good universal semanticity ranking. Important non-RI
writers and readers quantify its omissions. Conversely, low precision of the
59-head list does not justify saying that SIH has no connection to the mechanism.

\section{Frozen validation and stopping}
\label{sec:validation}
Freeze the code/model/input hashes, masks, selected exact routes, the enabled/disabled status of the optional query panel, one primary
conditional background per claimed member, reference means, metric signs, semantic
panel, thresholds and bootstrap procedure before opening the 87 held-out families.
Run one frozen suite: full model; retained $C$; empty head mask; complement/$C$
removal; RI-member removal; non-RI-member removal; the selected completeness groups;
and the frozen route/conditional tests supporting individual claims. Also evaluate the frozen reorder panel; include the four-condition panel only if enabled and frozen during discovery. The 87 families are new examples of the same
task, not different semantic relations. No selection or dropping of failures there.

Require the discovery faithfulness guards on validation and report intervals and
population-specific baseline accuracy. Mark each mechanistic claim confirmed,
inconclusive or not reproduced. One evaluation means a fixed suite, not one forward.
A post-validation revision is exploratory and needs fresh data for another
confirmatory claim.

Stop search after two expansion rounds or its allocated configuration limits;
report remaining gaps. If faithfulness never passes, validate the bounded mechanism
and its failure transparently rather than silently expanding to the whole model or
claiming a minimal circuit. Completion means the declared scope and limits have
been executed and reported, not that every hypothesis succeeds.

\section{Conditional extras, cost and deliverables}
\textbf{L30H18 redirect, only if its negative route remains unexplained.} Do not assume
its changed-query pattern places mass on the desired mother. Explicitly move its
colon attention mass on all mother spans onto the query mother's token span,
uniformly within that span, leaving all other mass and V unchanged. Compare with
moving the same mass to the other candidate mother and to a third mother. Save row
sums, intervention norms, both logits and candidate identities. Run the intact and
replaced head pattern on 40 pairs; extend only an informative result. This tests
the causal effect of routing under the fixed values, not the complete claim that
the head implements distractor copying. A changed-query donor can be a separate
natural-pattern control only after its actual destination is verified.

\textbf{Deferred.} No default SVD, new readout sweep, final-norm projection exercise,
parent-coverage test or new RI statistic. An actual matched residual-bypass test is
more informative than the old unmatched projection comparison. A revised L15H25
readout is warranted only if knowing its content would distinguish two surviving
mechanisms; first show source-swap sensitivity, then test head dependence. Broad
weight decompositions do not address the present decision before a route is known.

\textbf{Residual profiles.} During the necessary new baseline/group captures, save
norms and clean/corrupt distances after all blocks at mother-last, ``is'', child-last,
period and colon, including writer-restored differences where already computed.
These five-site profiles are inexpensive auxiliary outputs of those forwards, not
a separate search stage or a rule that excludes later heads. No old full-residual
cache is assumed.

\textbf{Cost accounting.} G1 has no exhaustive 25-head subset sweep. Its common-pair
cap is 128 configurations times 40 pairs = 5,120 intervention evaluations, a ceiling
rather than a required workload. G2 has 1,240 new conditional single evaluations,
plus at most $7\times2\times40=560$ initial group/reference evaluations and backgrounds.
G3 adds only nonduplicate configurations. Route capture and final injection are
separate costs; share caches only for identical frozen/live specifications.
Retain the cap of 96 full-discovery route contrasts in two directions:
$96\times178\times2=34,176$ endpoint evaluations, excluding capture passes.
The actual job manifest must count unique configurations and selected extensions.
Load the model once per worker, batch compatible inputs and reuse verified caches;
each distinct intervention still requires fresh affected computation. Do not quote
a runtime until the pilot measures it. The previous 3.7-hour figure was wall time,
not a transferable GPU-hour estimate.

Pilot six pairs and representative group, hybrid-route and optional-gradient jobs
on one two-GPU replica; record peak memory and seconds separately. Then estimate
wall time as the actual per-type work divided by independently allocated replicas,
plus load/cache/IO overhead. A six-GPU allocation supports three such replicas, not
three-way model parallelism by assumption. Use the established cluster submission
workflow; this document contains no submitted jobs or claimed runtime result.

Store prompt/model/code hashes, donor/recipient IDs, source and receiver masks,
channels and frozen/live branches, precomputed mean provenance, all logits,
attention/channel diagnostics, no-op drift and completion checksums for every
configuration. Resume only matching manifests, and verify exact ID coverage, not
only record counts. Suggested outputs: \file{group_events},
\file{route_events}, \file{receiver_blocking}, \file{family_summary},
\file{circuit_masks}, \file{faithfulness}, \file{ri_membership},
\file{search_boundary} and a separate immutable validation export.

The execution order is: prepare/gate and reuse or collect necessary baselines;
coverage and staged P mapping; route-supported G1 with receiver blocking; focused
G2 and G3 (which may run alongside mapping); triggered S4.3 expansion and return;
S4.4 order robustness/optional query panel and S4.5 reduction/completeness;
freeze; one held-out suite. No Stage-4 GPU experiment is recorded by this document.

\begin{thebibliography}{9}
\bibitem{ioi} K. Wang et al. (2022). \emph{Interpretability in the Wild: a Circuit
for Indirect Object Identification in GPT-2 small}.
\url{https://arxiv.org/abs/2211.00593}. Inspiration for exact paths, backups and
faithfulness/completeness/minimality; the thresholds here are project-specific.
\bibitem{patch} S. Heimersheim and N. Nanda (2024). \emph{How to Use and Interpret
Activation Patching}. \url{https://arxiv.org/abs/2404.15255}.
\bibitem{best} F. Zhang and N. Nanda (2024). \emph{Towards Best Practices of
Activation Patching in Language Models: Metrics and Methods}.
\url{https://arxiv.org/abs/2309.16042}.
\bibitem{eap} A. Syed, C. Rager and A. Conmy (2023). \emph{Attribution Patching
Outperforms Automated Circuit Discovery}. \url{https://arxiv.org/abs/2310.10348}.
\bibitem{faith} M. Hanna, S. Pezzelle and Y. Belinkov (2024). \emph{Have Faith in
Faithfulness: Going Beyond Circuit Overlap When Finding Model Mechanisms}.
\url{https://arxiv.org/abs/2403.17806}.
\bibitem{robustness} J. Miller, B. Chughtai and W. Saunders (2024).
\emph{Transformer Circuit Faithfulness Metrics are not Robust}.
\url{https://arxiv.org/abs/2407.08734}.
\end{thebibliography}
\end{document}
